\section{城市房地产管理法}
\begin{enumerate}
  \item 大量重要考点。\bs{开、转、变、售、交}。
  \item 房地产开发。
        \begin{enumerate}
          \item 用地。出让、划拨。
          \item 限期动工。防止囤地、捂地。
          \item 1 年+，\zd{闲置费 \xydy 出让金 20\%}。2 年未动工，无偿收回。
        \end{enumerate}
  \item 建设用地转让。\zd{房地一体}。
        \begin{enumerate}
          \item 出让地转让。\bs{2 金 2 证}。\zd{交全}出让金、完成 \zd{25\%} 的投资；办了土地使用证、Opt(房屋所有权证)。
          \item 划拨地转让。\yi 给民营企业。\er 给国家机关。
        \end{enumerate}
  \item 变更土地用途。\bs{2 同意 1 补签}。原出让方、城乡规划部门。变更协议 \hz 重新签订\zd{土地使用权出让合同}。
  \item 商品房的预售。\bs{3 证 2 金}。
        \begin{enumerate}
          \item 土地使用证、建设工程规划许可证、预售许可证。
          \item 交全出让金、完成 \bs{25\%} 投资。
          \item 预售合同备案（\zd{网签}）。不是强制义务（行政报批义务）。\scx{不是当事人权利}。
        \end{enumerate}
  \item 房地产交易。\zd{认地不认人}。不管怎么变，\bs{土地使用年限都不会重新计算}。
  \item 房地产抵押登记找县以上政府规定的部门办理。
  \item 房地产中介。咨询、价格评估、经纪机构。
\end{enumerate}
